% THETA3-B26 contribution by /root/definite_bridge, 21 November 2026.
\subsection{THETA3-B26: the mixed source Gram reduces to algebraic Hecke blocks and their multiplicity periods}\label{B26}

Keep the actual fixed kernel of B17 and B25.  There is a further
reduction of its source Gram matrix which does not replace the
kernel by a finite frame.  First its coefficient functions are
holomorphic sections of a specified algebraic automorphic line,
and their coherent coordinates have one common comparison scalar.
Second the actual Petersson form is constrained by an algebraic
Hecke algebra with adjoints.  On one irreducible Hecke module this
already makes the actual Gram projectively algebraic.  In general
the remaining quantities are positive Hermitian forms on the
Hecke multiplicity spaces.  The cohomological calculation below
identifies the periods represented by those remaining forms.

\paragraph{The actual source line and its central factor.}
Let
\[
 I=\{\tau\ne\mu:d_1(\tau)\ne d_2(\tau)\},\qquad s=|I|.
 \tag{B26a}
\]
The source $S=\U(P)$ has signature $(1,1)$ at the places in $I$
and is definite at the other real places.  Its adelic quotient
is compact because it is definite at $\mu$.  We discuss $s>0$;
the case $s=0$ is B24.  Choose the complex orientation of each
disk to be the one in which the actual Gaussian is holomorphic.
The opposite oscillator convention conjugates that disk and
every occurrence of its holomorphic line; it is not silently
identified with the chosen orientation.

At a place in $I$ the target $V_\tau$ is positive definite of
rank three.  On the source derived group the Gaussian has
lowest compact weight three.  Here is the direct Fock check.
For the compact torus of $\SU(1,1)$, the two source lines have
weights $e^{i\theta}$ and $e^{-i\theta}$.  In the positive
Lagrangian specified in 032, Section 4.6, printed p.22, the
second line is dualized.  Thus all six complex coordinates have
weight $e^{i\theta}$, and the half-determinant vacuum line has
weight $e^{3i\theta}$.  A unitary splitting changes only the
determinant character and is trivial on this derived torus.
The Gaussian is killed by the lowering root.  Equivalently,
the compact rank-three input with shifted parameter $(1,0,-1)$
has source parameter $(1;-1)+n_0/2$ in B18; its difference is
two and its holomorphic discrete-series weight is three.
The actual coefficient-vector convention may dualize both
weights, which is exactly the orientation convention just fixed.

At a definite source place different from $\mu$ the Gaussian
transforms by a determinant character.  At $\mu$ the degree-two
cochain uses the exterior square of the multiplicity factor;
this is again a determinant character of $\U(2)$.  Consequently
the complete source compact type is scalar.  Let $Y$ be its
compact unitary Shimura variety at a sufficiently small level,
with all the adelic components and the fixed central character
retained.  One may use a PEL similitude realization and then its
unitary components.  On a deeper principal component the
arithmetic determinant is one: the determinants of integral
unitary elements are norm-one algebraic units, hence roots of
unity, and a sufficiently small principal level kills them.
This does not remove their Hodge-theoretic determinant lines.

Let $\lambda_\tau$ denote the algebraic Hodge line whose
automorphy factor on the oriented upper half-plane is
$j(\gamma,z)=cz+d$.  The source line is
\[
 \mathscr L=
       \bigotimes_{\tau\in I}\lambda_\tau^{\otimes3}
                  \otimes\mathscr C.
 \tag{B26b}
\]
This equation defines $\mathscr C$ with the actual splitting
and compact characters: it is the quotient of the actual
one-dimensional compact type by the displayed weight-three
characters.  At each real place its remaining character is
an integral determinant power.  Indeed the full unitary
weights are integral after the specified unitary splitting,
and their difference at a switch is three.  Thus the quotient
is an algebraic character of the determinant torus.  Its
finite character part is retained as well.

The bundle $\mathscr C$ is the corresponding fixed CM
determinant line, with a possible finite character local
system.  On the finite component cover just described its
flat monodromy is trivial.  Its Hodge filtration and comparison
period need not be trivial.  Choose a unit flat analytic frame
$u_0$ and an algebraic de Rham frame $u_{\rm dR}$ of this line,
and record
\[
                 u_{\rm dR}=\Omega_{\mathscr C}u_0,
             \qquad \|u_{\rm dR}\|^2=|\Omega_{\mathscr C}|^2.
 \tag{B26c}
\]
If several labelled components are used, transport the frames
by their specified finite Hecke isogenies; their algebraic
transport factors are retained in the component coordinates.
Nothing sets $\Omega_{\mathscr C}$ to one or declares the line
to have Hodge weight zero.

The algebraicity of (B26b), its finite Hecke action and its
comparison at CM fibres can be read directly from the PEL
Hodge bundles.  In general canonical-bundle language the
precise source is Milne, \emph{Canonical models of mixed
Shimura varieties and automorphic vector bundles}, corrected
2018 version, Chapter III, Remark 2.3 and Example 2.4,
pp.55--56; Theorem 5.1 with proof and Remark 5.2, p.61;
Section 6, pp.61--62; and Section 7, pp.62--63.
These give algebraic automorphic bundles and their
$S(\A_{F,f})$-equivariant structures, not rational Betti
coordinates of arbitrary holomorphic sections.

Write the actual kernel in B25's algebraic coherent target
basis as
\[
                  K(g,h)=\sum_{a=1}^r f_a(g)e_a(h).
 \tag{B26d}
\]
The lowering-root equation for the source Gaussian survives
application of each target coefficient functional.  Therefore
each $e_a$ is the scalar automorphic realization of a
holomorphic section of $\mathscr L$.  This is a finite-dimensional
space because $Y$ is projective.  No finite-dimensionality
claim for arbitrary fixed-compact-type smooth functions is used.

\paragraph{One common coherent scalar on the actual source.}
Choose the source CM point $y$ defined by the original rational
decomposition $P=P_1\perp P_2$.  Its Shimura cocharacter lies
in the product of the two line tori.  At a switch it selects
the indicated positive or negative line; at a definite place
it has the prescribed central lift.  This establishes the CM
point for the full unitary datum, not merely a map from its
derived subgroup.

The finite adelic translates of $y$ separate
$H^0(Y,\mathscr L)$.  To see this, represent a section by its
scalar automorphic function $a(h_\infty,h_f)$.  If this is
nonzero at some point, choose $\gamma\in S(F)$ whose real
component is sufficiently close to $h_\infty$.  The Cayley
weak approximation argument of B25 gives such $\gamma$, and
automorphy gives
\[
       a(1,\gamma_f^{-1}h_f)=a(\gamma_\infty,h_f)\ne0.
 \tag{B26e}
\]
Because the section space is finite-dimensional, finitely
many such evaluations separate it.

Fix an algebraic basis $\eta_1,\ldots,\eta_D$ of
$H^0(Y,\mathscr L)$ and finitely many separating translates
$b_l$.  Transport an algebraic fibre frame from $y$ to the
CM point $[y,b_l]$ by the specified finite isogeny.  In the
marked analytic uniformization that isogeny has the identity
archimedean coordinate.  Hence the conversion from these
algebraic frames to the fixed automorphic frame has one
common nonzero scalar $t_y$, with any algebraic isogeny
factors put in the matrix.  The evaluation matrix is
\[
                   R_y=t_y B,
           \qquad B\in M_{L,D}(\overline\Q),\quad
                       \operatorname{rank}B=D.
 \tag{B26f}
\]
This is the same fibre-transport argument as B23, applied to
the actual line (B26b); it does not use CM of the cohomology
of $Y$.  It also retains all periods of $\mathscr C$ in $t_y$.

Include these finitely many source translates in the common
finite theta family of B19 and B25.  Assume its specified
finite Schwartz and splitting values are algebraic, exactly
as in B25.  The coefficient vector of
$K(\,\cdot\,,(1,b_l))$ in the basis $f_a$ is
$(C_x/t_x)$ times an algebraic vector.  Here $C_x$ is the
common CM theta comparison for this combined finite family,
and $t_x$ is the target coherent-fibre scalar.  Enlarging the
finite theta level is done once for this family, with its
isogeny factors retained.  Thus, for some algebraic matrix $D_0$,
\[
                    (e_a(1,b_l))_{l,a}
                         ={C_x\over t_x}D_0.
 \tag{B26g}
\]
Let $C$ be the coordinate matrix of the $e_a$ in the basis
$\eta$.  Taking an algebraic left inverse $L_B$ in (B26f)
gives the exact conclusion
\[
              C=\gamma C_0,\qquad
              \gamma={C_x\over t_x t_y},\qquad
                    C_0=L_BD_0\in M(\overline\Q).
 \tag{B26h}
\]
In particular the actual source coefficient space is defined
over $\overline\Q$ in coherent coordinates.  All its vectors
have a common scalar in the chosen coordinate list.  Neither
this statement nor Milne's CM-value criterion identifies
those coherent coordinates with Betti coordinates.

\paragraph{The exact cohomological source Gram.}
On each oriented disk use $z_\tau=x_\tau+iy_\tau$ and
$d\mu_\tau=dx_\tau dy_\tau/y_\tau^2$.  After removing the
constant line specified below, the local system is
\[
                     \mathbb V=
               \bigotimes_{\tau\in I}\mathbb V_\tau,
 \tag{B26i}
\]
where $\mathbb V_\tau$ is the standard rank-two weight-one
variation, in the complex embedding selected by the disk.
It is the corresponding factor of the first cohomology of
the source PEL abelian scheme.  The tensor can be made rational
by including its Galois conjugates; the displayed complex
summand is retained throughout the following calculation.
It has weight $s$.

Here the purity assertion has an explicit geometric
realization.  Let $\mathcal B\to Y$ be the source PEL
abelian scheme and $Z=\mathcal B\times_Y\cdots\times_Y
\mathcal B$ its $s$-fold fibre power.  Both $Y$ and $Z$
are smooth projective.  The indicated tensor local
system is a coefficient-field summand of
$R^s\pi_*\Q$, obtained by taking $R^1$ from each
factor and the prescribed $E$-embedding idempotents.
For an integer $N>1$, simultaneous multiplication
by $N$ on all fibres acts by $N^q$ on $R^q\pi_*\Q$.
All differentials in its Leray spectral sequence
therefore vanish: a differential of page $r\ge2$
would connect the distinct eigenvalues $N^q$ and
$N^{q-r+1}$.  In degree $2s$, the successive Leray
grades have distinct eigenvalues.  Hence a rational
polynomial in $[N]^*$ projects $H^{2s}(Z,\Q)$ onto
the grade $H^s(Y,R^s\pi_*\Q)$.  Indeed the product
of the distinct linear factors annihilates the
filtered vector space, so this operator is
semisimple and the usual Lagrange interpolation
polynomial is a projector.

Within that grade, multiplication by $N$ separately
on each fibre factor selects its $R^1$ summand by
the same interpolation argument, now on the
canonical fibre K\"unneth decomposition.  Finally
the $E$-endomorphism projectors select $\mathbb V$.
All these operators are algebraic graph
correspondences with the stated coefficient field.
Including the conjugate embedding projectors gives
the corresponding rational summand.  Thus the
cohomology used here is a prescribed summand of the
pure weight-$2s$ Hodge structure of $H^{2s}(Z)$.
The exact primary purity and polarization locators
are Deligne, \emph{Th\'eorie de Hodge II},
Section 2.2.1--2.2.6, printed pp.27--29.
Theorem 4.1.1(i), printed p.40, supplies the general
proper-smooth Leray degeneration statement; the
multiplication argument above proves the particular
splitting needed here directly.

There is a necessary determinant correction in this
cohomological interpretation.  The tangent bundle of the
line domain is $\operatorname{Hom}(\lambda_\tau,
\mathbb V_\tau/\lambda_\tau)$, so its canonical line is
$\lambda_\tau^2\otimes(\det\mathbb V_\tau)^{-1}$.
Consequently
\[
 \mathscr D=\mathscr C\otimes
                 \bigotimes_{\tau\in I}\det\mathbb V_\tau,
 \qquad
 \mathscr L=\Omega_Y^s\otimes F^s\mathbb V\otimes\mathscr D.
 \tag{B26i$'$}
\]
Here the second equality is the equivariant
Kodaira--Spencer identification, with its fixed
normalization.  The extra determinants are constant
CM Hodge lines.  Their selected embeddings have type
$(1,1)$, but their other conjugates need not make them
globally Tate.  Trivial determinant monodromy does not
discard these factors.  Use the product de Rham frame
of $\mathscr D$ and record its comparison with a unit
flat frame as $\Omega_{\mathscr D}$.  Thus
$\Omega_{\mathscr D}=\Omega_{\mathscr C}
\prod_{\tau\in I}\Omega_{\det\mathbb V_\tau}$ in the
chosen product frames, with every specified
Kodaira--Spencer frame factor retained.  In particular
any universal $2\pi i$ factors in these determinant
comparisons are not set equal to one.

On a component, write the actual coefficient in the form
\[
           e_a(g_z)=
              \Omega_{\mathscr D}F_a(z)
                          \prod_{\tau\in I}y_\tau^{3/2}
 \tag{B26j}
\]
in the fixed unitary automorphic frame.  This defines $F_a$;
no assertion that $F_a$ has algebraic Betti periods is being
made.  Put $v_\tau(z_\tau)=(z_\tau,1)$ and
\[
 \alpha_a=F_a(z)\,
       \bigotimes_{\tau\in I}v_\tau(z_\tau)
                   \;dz_1\wedge\cdots\wedge dz_s.
 \tag{B26k}
\]
For $\gamma_\tau\in\mathrm{SL}_2(\R)$ one has
$v_\tau(\gamma_\tau z)=
 \gamma_\tau v_\tau(z)/j(\gamma_\tau,z)$ and
$d(\gamma_\tau z)=dz/j(\gamma_\tau,z)^2$.
The weight-three transformation of $F_a$ therefore makes
(B26k) a well-defined $\mathbb V$-valued form.  It is closed
for the flat connection: its coefficients are holomorphic,
and every holomorphic derivative wedges with the already
present top product of the $dz_\tau$.
Its class has type $(2s,0)$ in the weight-$2s$ cohomology
$H^s(Y,\mathbb V)$, after the stated untwisting of
$\mathscr D$.  Restoring that line restores its actual CM
Hodge type and period; it is not a tacit Tate twist.

For the flat alternating form
$\varepsilon((a,b),(c,d))=ad-bc$ let
$B_{\mathbb V}=\bigotimes_{\tau\in I}\varepsilon_\tau$.
Define the Hermitian pairing, conjugate linear in the first
variable, on these classes by
\[
 Q(\alpha_a,\alpha_b)=
  (-1)^{s(s-1)/2}
     \int_Y B_{\mathbb V}(\alpha_b\wedge\overline{\alpha_a}).
 \tag{B26l}
\]
The convention in this display pairs coefficient tensors and
wedges the displayed base forms in that order.  Since
$\varepsilon(v,\bar v)=2iy$ and
$dz\wedge d\bar z=-2i\,dx\wedge dy$, moving the conjugate
base factors next to their partners gives exactly
\[
 Q(\alpha_a,\alpha_b)=
       4^s\int_Y\overline{F_a(z)}F_b(z)
                     \prod_{\tau\in I}y_\tau^3d\mu_\tau.
 \tag{B26m}
\]
In particular this form is positive definite on the
coefficient classes.  If a nonzero $\alpha_a$ were exact,
Stokes's theorem paired with its closed conjugate would make
the left side zero, contradicting the positive right side.
Thus no nonzero source coefficient is lost on passage to
this cohomology.

Let $\nu_S>0$ be the conversion from the specified adelic
Haar measure to the displayed disk measures and the fixed
compact-fibre and component measures.  Keep it as the
specified measure conversion, rather than renormalizing the
source quotient.  At one fixed source level, disintegration
gives the same finite-level volume and compact-fibre volume
on every adelic component.  If a deeper principal component
cover was used above, its contribution is divided by its
integer covering degree.  Thus the full disjoint-union
pairing in (B26l) includes these specified rational component
weights, and a single common positive conversion remains
in $\nu_S$.  The actual
conjugate-linear-first source Gram is
\[
 \widehat W_{ab}:=\int_{[S]}\overline{e_a(h)}e_b(h)\,dh
      ={\nu_S|\Omega_{\mathscr D}|^2\over4^s}
                                      Q(\alpha_a,\alpha_b).
 \tag{B26n}
\]
B17 uses $W_{ab}=\int e_a\overline{e_b}\,dh$ for its
coefficient matrix, so $W=\widehat W^{\mathsf t}$.
This transpose is retained when (B26n) is substituted in
the actual $TT^*$ matrix.

This gives a concrete possible CM criterion.  If the classes
of the normalized coherent source sections lie in one fixed
coefficient-embedding block of a full simple CM isotypic
summand, their algebraic de Rham classes
have a common Betti comparison period, just as in B23.
Their matrix in (B26n) is then one positive real scalar
times an algebraic Hermitian matrix.  Such a CM statement
is a statement about $H^s(Y,\mathbb V)$ with its recorded
constant CM line.  It does not follow from CM of the target
cup image.  In particular there has been no unproved
rational cohomological theta correspondence from the target
to this source group.

\paragraph{Algebraic Hecke adjoints reduce the remaining matrix.}
There is a separate reduction that needs no source CM
hypothesis.  Let $\mathcal H_{\overline\Q}$ be the finite-level
source Hecke algebra acting on $H^0(Y,\mathscr L)$, with the
actual central character imposed.  Normalize the compact
open subgroup to have finite Haar volume one in the
convolution notation.  Double-coset sums and their degrees
then have rational coefficients.  This convention does not
change the already fixed quotient measure in $\widehat W$.
The geometric Hecke correspondences and their bundle maps
are defined over $\overline\Q$, by the equivariant canonical
model above.  Thus they act algebraically in the coherent
basis.

The group is unimodular.  Change of variables in the
actual Petersson integral proves
\[
    \langle T_\phi u,v\rangle_S
        =\langle u,T_{\phi^*}v\rangle_S,
       \qquad \phi^*(h)=\overline{\phi(h^{-1})}.
 \tag{B26o}
\]
For characteristic functions this is the inverse double
coset, with the same convolution normalization.  Hence
the algebraic Hecke image is closed under the adjoint
involution of the actual positive Hermitian form.

Let $\mathcal F_{\overline\Q}$ be the Hecke saturation of
the normalized sections $\gamma^{-1}e_a$ from (B26h).
It is a finite-dimensional algebraic subspace of
$H^0(Y,\mathscr L)$.  Let $\mathcal A$ be the image of the
Hecke algebra on it.  Its complexification is a finite-dimensional
algebra of operators closed under the Hermitian adjoint.
Every invariant subspace therefore has an invariant
orthogonal complement.  Thus it is semisimple, and so is
$\mathcal A$ over $\overline\Q$.  Choose an algebraic
Wedderburn decomposition of this actual representation:
\[
 \mathcal F_{\overline\Q}
    =\bigoplus_\lambda M_\lambda\otimes V_\lambda,
 \qquad
 \mathcal A=\bigoplus_\lambda\operatorname{End}(V_\lambda),
 \qquad m_\lambda=\dim M_\lambda.
 \tag{B26p}
\]
The central block idempotents are orthogonal projections
for the actual Hermitian form.  This follows either from
the invariant-complement assertion or from the fact that
a central idempotent in a finite-dimensional adjoint-closed
operator algebra is self-adjoint.

For each $V_\lambda$ there is a positive definite algebraic
Hermitian matrix $G_\lambda$ implementing the given
involution.  To justify algebraicity, choose a finite
algebraic spanning set of the matrix algebra.  The equations
\[
        G_\lambda\rho(a)=\rho(a^*)^*G_\lambda
 \tag{B26q}
\]
are a finite real linear system with real algebraic
coefficients on the Hermitian entries.  A positive real
solution exists by restriction of the actual form to an
irreducible copy.  Algebraic points are dense in this
linear solution space, and positivity is open, so a
positive algebraic solution exists.  Algebraic changes
of basis, including an algebraic Cholesky factor, are
therefore allowed here.

Schur's lemma now gives the exact shape of the actual
Petersson Gram in (B26p):
\[
                   G_S=\bigoplus_\lambda
                                  H_\lambda\otimes G_\lambda,
 \tag{B26r}
\]
where each $H_\lambda$ is a positive definite Hermitian
matrix on the multiplicity space $M_\lambda$.  These
$H_\lambda$ are actual source period matrices; no
algebraicity is asserted for them.  For completeness,
compose the actual form on one block with the inverse
of $I\otimes G_\lambda$.  Invariance makes the resulting
operator commute with $I\otimes\operatorname{End}(V_\lambda)$,
so it has the form $H_\lambda\otimes I$.  Positivity and
Hermitian symmetry give the stated properties of
$H_\lambda$.

Let $C_1$ be the algebraic matrix of the normalized
sections $\gamma^{-1}e_a$ in (B26p).  Their actual Gram
is therefore
\[
       \widehat W=|\gamma|^2 C_1^*
          \left(\bigoplus_\lambda H_\lambda\otimes G_\lambda\right)
                                  C_1.
 \tag{B26s}
\]
This is a reduction of the Gram of the original fixed
kernel, not a replacement by norms of a different frame.

If its Hecke saturation is one irreducible module, then
there is one index $\lambda$ and $m_\lambda=1$.
Equation (B26s) becomes
\[
                 \widehat W=c\,C_1^*G_\lambda C_1,
                         \qquad c>0.
 \tag{B26t}
\]
The same conclusion holds if the coefficient vectors
occupy a single algebraic line in a single multiplicity
space.  In these cases the actual $TT^*$ inverse of B17
is one positive scalar times an algebraic matrix on
the preserved algebraic image of B25.  B23 realizes that
matrix by an algebraic CM-source correspondence.  Thus
the actual fixed-kernel correction is closed in this
stated Hecke sector, without proving the source motive
to be CM.

If several multiplicity-one sectors occur, each has its
own positive scalar; their ratios have not been shown
algebraic.  If $m_\lambda>1$, the remaining matrix has
$m_\lambda^2$ real Hermitian parameters before any further
relations.  The CM criterion after (B26n) is one way to
control these quantities, but it is an additional
source-motive assertion.

\paragraph{A finite sufficient certificate.}
The irreducible-sector condition can be certified without
assuming a global multiplicity-one theorem.  It suffices
to exhibit a finite-dimensional algebraic coherent
subspace containing all $\gamma^{-1}e_a$ and a finite
list of algebraic geometric Hecke operators and their
adjoints preserving it, whose matrices generate its full
matrix algebra.  Generation is a finite linear-algebra
test on products of these matrices; once their span has
dimension $D^2$, it is the full matrix algebra on a
$D$-dimensional space.  The proof of (B26q)--(B26t)
then applies using just that finite list.  The necessary
coherent matrices can be recovered from finite CM-value
rows with the same common-scalar argument (B26f)--(B26h),
retaining the original finite theta and isogeny data.
No such full-matrix certificate is claimed here for an
unspecified actual 032 input.

\paragraph{Source scope.}
The Fock and real representation inputs are the exact
032 Section 4.6 positive Lagrangians and B18's checked
Atobe Definition 1.6/Theorem 1.7 and Ichino Theorem 4.1
specialization.  The finite CM-value comparison is B19
and B25 with its original theta-group, tangent, splitting
and Haar normalizations.  Milne's Chapter III supplies
the algebraic coherent bundle and Hecke structures.
The cohomological factor $4^s$, its sign and the entire
finite-dimensional adjoint-algebra reduction have been
proved here.  None of these statements identifies a sum
of sectors or a changed finite kernel with one prescribed
arithmetic \texttt{MainX} tuple.
